% ------------------------------------------------------------------------------
% Chapter 1
% ------------------------------------------------------------------------------
\chapter{Design Overview} 
\label{cha:design_overview}

There are 2 modules under test: fifo and cache.

\section{Fifo}

A circular buffer.
It utilizes internal write and read pointers to manage data flow and provides status flags to prevent data loss or invalid reads.

\subsection{Parameters}

There are two parameters to this dut:

\begin{itemize}
    \item \texttt{WIDTH} : data size.
    \item \texttt{DEPTH} : fifo depth.
\end{itemize}

\subsection{Interface}

\begin{table}[h!]
    \centering
    \begin{tabular}{|l|c|c|l|}
        \hline
        \textbf{Signal}     & \textbf{Direction} & \textbf{Width} & \textbf{Description}            \\ \hline
        \texttt{clk}        & Input              & 1              & Clock signal                    \\ \hline
        \texttt{rst\_n}     & Input              & 1              & Active-low synchronous reset    \\ \hline
        \texttt{wr\_en}     & Input              & 1              & Write enable signal             \\ \hline
        \texttt{rd\_en}     & Input              & 1              & Read enable signal              \\ \hline
        \texttt{wr\_data}   & Input              & WIDTH          & Write data                      \\ \hline
        \texttt{rd\_data}   & Output             & WIDTH          & Read data                       \\ \hline
        \texttt{full}       & Output             & 1              & Full signal                     \\ \hline
        \texttt{empty}      & Output             & 1              & Empty signal                    \\ \hline
    \end{tabular}
    \caption{Fifo signals}

\end{table}

\subsection{Functional Description}

\begin{itemize}

    \item \textbf{Reset State}: On a low-level reset (\texttt{rst\_n}), all internal pointers and the occupancy count are cleared. The empty flag is asserted, and the full flag is de-asserted.

    \item \textbf{Write Operation}: Triggered by \texttt{wr\_en}. If the FIFO is not full and no simultaneous read is requested, data is written to the current \texttt{wr\_ptr} location, the pointer increments (modulo \texttt{DEPTH}), and the count increases.
    
    \item \textbf{Read Operation}: Triggered by \texttt{rd\_en}. If the FIFO is not empty and no simultaneous write is requested, data is retrieved from \texttt{rd\_ptr}, the pointer increments (modulo \texttt{DEPTH}), and the count decreases.

    \item \textbf{Conflict Resolution}: The logic includes a strict mutual exclusion rule: if both \texttt{wr\_en} and \texttt{rd\_en} are high, the FIFO performs no action, maintaining the current state.

    \item \textbf{Flag Logic}:
    \begin{itemize}
        \item \textbf{Full}: Asserted when count == \texttt{DEPTH}.
        \item \textbf{Empty}: Asserted when count == 0.
    \end{itemize}

\end{itemize}

\section{Cache}

A direct-mapped cache.
It translates an 8-bit address into Tag, Index, and Offset fields to determine if requested data is currently "local" (hit) or needs to be fetched (miss).

\subsection{Parameters}

There are four parameters to this dut:

\begin{itemize}
    \item \texttt{CACHE\_LINES} : lines in the cache.
    \item \texttt{LINE\_SIZE} : size of each line.
    \item \texttt{ADDR\_WIDTH} : size of the address.
    \item \texttt{DATA\_WIDTH} : size of the data.
\end{itemize}

\subsection{Interface}

\begin{table}[h!]
    \centering
    \begin{tabular}{|l|c|c|l|}
        \hline
        \textbf{Signal}     & \textbf{Direction} & \textbf{Width} & \textbf{Description}    \\ \hline
        \texttt{clk}        & Input              & 1              & Clock signal            \\ \hline
        \texttt{reset}      & Input              & 1              & synchronous reset       \\ \hline
        \texttt{read}       & Input              & 1              & read signal             \\ \hline
        \texttt{write}      & Input              & 1              & write signal            \\ \hline
        \texttt{addr}       & Input              & ADDR\_WIDTH    & address                 \\ \hline
        \texttt{data\_in}   & Input              & DATA\_WIDTH    & input data              \\ \hline
        \texttt{data\_out}  & Output             & DATA\_WIDTH    & output data             \\ \hline
        \texttt{hit}        & Output             & 1              & hit signal              \\ \hline
    \end{tabular}
    \caption{Cache signals}

\end{table}

\subsection{Functional Description}

\begin{itemize}

    \item \textbf{Address Decoding}: The input address (\texttt{addr}) is partitioned into three functional fields:
    \begin{itemize}
        \item \textbf{Tag}: Identifies the memory block.
        \item \textbf{Index}: Selects one of the 16 cache lines.
        \item \textbf{Offset}: Selects the specific 32-bit word within the 4-word line.
    \end{itemize}

    \item \textbf{Read Operation (Priority)}: 
    \begin{itemize}
        \item \textbf{Hit}: If the indexed line is valid and the stored tag matches the request, the \texttt{hit} signal is asserted and \texttt{data\_out} is driven from internal memory.
        \item \textbf{Miss}: If the tag differs or the line is invalid, \texttt{hit} is de-asserted. The cache performs a simulated fetch, updating the \texttt{cache\_mem}, \texttt{tag\_array}, and \texttt{valid\_array} with the new data.
    \end{itemize}
    
    \item \textbf{Write Operation}: Updates the \texttt{cache\_mem} with \texttt{data\_in}. 
    \begin{itemize}
        \item If the access matches an existing valid tag, it is reported as a \texttt{hit}.
        \item If the access is a \texttt{miss}, the \texttt{tag\_array} and \texttt{valid\_array} are updated to reflect the new memory block.
    \end{itemize}

    \item \textbf{Precedence Rule}: The module implements a strict priority for read requests. If both \texttt{read} and \texttt{write} signals are asserted simultaneously, the design executes the read operation only, and the write request is ignored.

    \item \textbf{Replacement Policy}: Being a direct-mapped cache, it utilizes an implicit "always overwrite" policy. Any access to an occupied index with a non-matching tag results in the previous data being replaced without further logic.

\end{itemize}